\documentclass[11pt]{article}
\usepackage[margin=1in]{geometry}
\usepackage{amsmath,amssymb,amsthm}
\usepackage{natbib}
\usepackage{hyperref}
\usepackage{booktabs}
\usepackage{enumitem}

\newcommand{\E}{\mathbb{E}}
\newcommand{\R}{\mathbb{R}}
\newcommand{\bth}{\boldsymbol{\theta}}
\newcommand{\bpsi}{\boldsymbol{\psi}}
\newcommand{\by}{\mathbf{y}}
\newcommand{\cM}{\mathcal{M}}
\newcommand{\ptilde}{\tilde{p}}

\theoremstyle{plain}
\newtheorem{proposition}{Proposition}
\theoremstyle{definition}
\newtheorem{definition}{Definition}
\newtheorem{example}{Example}
\newtheorem{remark}{Remark}

\title{Bridge Sampling for Marginal Likelihood Estimation:\\
Theory, Implementation, and Diagnostics for Hierarchical RDU Models}
\author{Technical Note}
\date{May 2026}

\begin{document}
\maketitle

\tableofcontents
\newpage

\section{Problem Statement}

Given a model $\cM$ with parameters $\bth \in \R^d$, data $\by$, and prior
$\pi(\bth)$, the marginal likelihood is
\begin{equation}
  \label{eq:ml}
  p(\by \mid \cM) = \int p(\by \mid \bth, \cM)\, \pi(\bth \mid \cM)\, d\bth.
\end{equation}
Define the unnormalized posterior
$\ptilde(\bth) \equiv p(\by \mid \bth)\, \pi(\bth)$, so that
$p(\by \mid \cM) = \int \ptilde(\bth)\, d\bth$. The posterior is
$p(\bth \mid \by) = \ptilde(\bth) / p(\by \mid \cM)$.

MCMC gives us samples from $p(\bth \mid \by)$ and the ability to evaluate
$\ptilde(\bth)$ at any point, but not $p(\by \mid \cM)$ itself. Bridge
sampling estimates $p(\by \mid \cM)$ from these ingredients plus samples from
a proposal distribution.


\section{The Bridge Sampling Identity}

\subsection{Statement}

Let $q(\bth)$ be a normalized proposal distribution from which we can both
sample and evaluate the density. \citet{MengWong1996} show that for any
function $h(\bth) > 0$ satisfying integrability conditions,
\begin{equation}
  \label{eq:bridge_identity}
  p(\by \mid \cM)
  = \frac{\E_q\!\left[\ptilde(\bth)\, h(\bth)\right]}
         {\E_{\text{post}}\!\left[q(\bth)\, h(\bth)\right]}.
\end{equation}

\subsection{Derivation}

Write the numerator as
$\int \ptilde(\bth)\, h(\bth)\, q(\bth)\, d\bth$. Write the denominator as
\[
  \int q(\bth)\, h(\bth)\, p(\bth \mid \by)\, d\bth
  = \int q(\bth)\, h(\bth)\, \frac{\ptilde(\bth)}{p(\by \mid \cM)}\, d\bth
  = \frac{1}{p(\by \mid \cM)} \int \ptilde(\bth)\, h(\bth)\, q(\bth)\, d\bth.
\]
The ratio cancels the common integral, leaving $p(\by \mid \cM)$. The identity
holds for any positive $h$; the choice of $h$ affects only the variance of the
estimator, not its expectation.


\section{The Optimal Bridge Function and Iterative Estimator}

\subsection{Optimal $h$}

\citet{MengWong1996} show the variance-minimizing bridge function is
\begin{equation}
  \label{eq:optimal_h}
  h^*(\bth) = \frac{1}{s_1\, \ptilde(\bth) + s_2\, q(\bth)\, r}
\end{equation}
where $r = p(\by \mid \cM)$ is the unknown quantity,
$s_1 = N_1 / (N_1 + N_2)$, $s_2 = N_2 / (N_1 + N_2)$, with $N_1$ proposal
samples and $N_2$ posterior samples.

\subsection{The circularity and the iteration}

Since $h^*$ depends on $r$, the estimator is defined implicitly. The solution
is to iterate. Given posterior samples
$\{\bth^{(i)}_\text{post}\}_{i=1}^{N_2}$ and proposal samples
$\{\bth^{(j)}_q\}_{j=1}^{N_1}$:

\begin{enumerate}[label=\textbf{Step \arabic*.}, leftmargin=2.5cm]
  \item \textbf{Initialize} $\hat{r}^{(0)}$. A reasonable starting value is
    the median of $\ptilde(\bth^{(i)}_\text{post}) / q(\bth^{(i)}_\text{post})$
    over the posterior samples, since
    $\E_\text{post}[\ptilde / q] = r$ for a perfect proposal.

  \item \textbf{Compute the numerator} (average over proposal samples):
    \begin{equation}
      \label{eq:numerator}
      \ell_1^{(t)} = \frac{1}{N_1} \sum_{j=1}^{N_1}
        \frac{\ptilde(\bth^{(j)}_q)}
             {s_1\, \ptilde(\bth^{(j)}_q) + s_2\, q(\bth^{(j)}_q)\, \hat{r}^{(t)}}
    \end{equation}

  \item \textbf{Compute the denominator} (average over posterior samples):
    \begin{equation}
      \label{eq:denominator}
      \ell_2^{(t)} = \frac{1}{N_2} \sum_{i=1}^{N_2}
        \frac{q(\bth^{(i)}_\text{post})}
             {s_1\, \ptilde(\bth^{(i)}_\text{post}) + s_2\,
              q(\bth^{(i)}_\text{post})\, \hat{r}^{(t)}}
    \end{equation}

  \item \textbf{Update}: $\hat{r}^{(t+1)} = \ell_1^{(t)} / \ell_2^{(t)}$.

  \item \textbf{Repeat} until $|\hat{r}^{(t+1)} - \hat{r}^{(t)}| < \epsilon$.
\end{enumerate}

\subsection{Log-scale computation}

All quantities are computed on the log scale to avoid overflow. Define
$\text{LSE}$ as the log-sum-exp operation:
$\text{LSE}(a_1, \ldots, a_n) = \log \sum_i e^{a_i}$.

\begin{align}
  \log \hat{r}^{(t+1)} &= \log \ell_1^{(t)} - \log \ell_2^{(t)} \\
  \log \ell_1^{(t)} &= -\log N_1 + \text{LSE}_{j=1}^{N_1}\!\left[
    \log \ptilde_j - \text{LSE}\!\left(
      \log s_1 + \log \ptilde_j,\;
      \log s_2 + \log q_j + \log \hat{r}^{(t)}
    \right)
  \right] \label{eq:log_num} \\
  \log \ell_2^{(t)} &= -\log N_2 + \text{LSE}_{i=1}^{N_2}\!\left[
    \log q_i - \text{LSE}\!\left(
      \log s_1 + \log \ptilde_i,\;
      \log s_2 + \log q_i + \log \hat{r}^{(t)}
    \right)
  \right] \label{eq:log_den}
\end{align}

where $\ptilde_j = \ptilde(\bth^{(j)}_q)$, $q_j = q(\bth^{(j)}_q)$, etc.


\subsection{Self-correcting behavior}

The iteration is a contraction mapping. If $\hat{r}^{(t)}$ is too large,
the $s_2 q r^{(t)}$ term in the denominator of $h^*$ dominates, which
inflates the denominator average relative to the numerator, driving
$\hat{r}^{(t+1)}$ downward. Conversely, if $\hat{r}^{(t)}$ is too small,
$s_1 \ptilde$ dominates and the update drives $r$ upward.

At the fixed point $\hat{r} = r$, the bridge function
$h^*(\bth) = 1/(s_1 \ptilde + s_2 q r)$ is a harmonic mean of $\ptilde$
and $qr$, weighting each sample by how ``surprising'' it is relative to both
distributions.


\section{What the Bridge Estimator Needs}

To run the estimator, we need four ingredients:

\begin{enumerate}
  \item \textbf{Posterior samples}
    $\{\bth^{(i)}_\text{post}\}_{i=1}^{N_2}$: from MCMC. These are the
    chain states saved in history files.

  \item \textbf{Proposal samples}
    $\{\bth^{(j)}_q\}_{j=1}^{N_1}$: drawn from $q$. These are
    \textit{not} MCMC samples. They are independent draws from a
    distribution we choose (e.g., multivariate normal).

  \item \textbf{Evaluations of $\log \ptilde(\bth)$} at every sample (both
    posterior and proposal). This is the full unnormalized log-posterior:
    \begin{equation}
      \label{eq:full_lp}
      \log \ptilde(\bth, \bpsi) = \log \pi(\bpsi)
        + \sum_{i=1}^{n} \left[
          \log p(\bth_i \mid \bpsi) + \log |J_g(\bth_i)| + \log p(\by_i \mid \bth_i)
        \right]
    \end{equation}
    where $\bpsi$ are hyperparameters, $p(\bth_i \mid \bpsi)$ is the
    population prior, $|J_g(\bth_i)|$ is the Jacobian correction, and
    $p(\by_i \mid \bth_i)$ is the data likelihood for subject $i$.

  \item \textbf{Evaluations of $\log q(\bth)$} at every sample. This requires
    a proposal with a known, computable density.
\end{enumerate}

Items 3 and 4 are the expensive parts. Item 3 requires looping over all $n$
subjects to compute the full joint, which is computationally costly but
parallelizable on GPU. Item 4 depends on the proposal choice.


\section{Sources of NaN and Numerical Failures}

\subsection{The float32 clipping problem}

The data likelihood for each choice computes a choice probability via the
logistic sigmoid:
\begin{align}
  \Delta &= \frac{V_L - V_R}{\lambda \cdot (u_{\max} - u_{\min})} \\
  P(\text{choose } L) &= \sigma(\Delta) = \frac{1}{1 + e^{-\Delta}}
\end{align}
The log-likelihood contribution for one choice is
\begin{equation}
  \label{eq:choice_ll}
  \ell_c = y_c \log P + (1 - y_c) \log(1 - P)
\end{equation}
where $y_c \in \{0, 1\}$ is the observed choice.

When $|\Delta|$ is large (strong preference), $\sigma(\Delta)$ saturates to 0
or 1. In IEEE float32, the sigmoid returns \textit{exactly} 1.0 for
$\Delta \gtrsim 16.6$ and exactly 0.0 for $\Delta \lesssim -16.6$.

\textbf{The clipping failure.} To prevent $\log(0)$, the code clips
$P$ to $[\epsilon, 1 - \epsilon]$ with $\epsilon = 10^{-10}$.
In float32, $1 - 10^{-10}$ rounds to $1.0$ because float32 has only
$\sim 7$ significant digits ($\epsilon_{\text{machine}} \approx 6 \times
10^{-8}$). The clip is a no-op, and $P = 1.0$ passes through unchecked.

\textbf{The $0 \times (-\infty) = \text{NaN}$ problem.} When $y_c = 1$ (chose
Left) and $P = 1.0$:
\[
  \ell_c = 1 \cdot \log(1.0) + 0 \cdot \log(1.0 - 1.0) = 0 + 0 \cdot (-\infty) = \text{NaN}
\]
IEEE 754 defines $0 \times \pm\infty = \text{NaN}$. Even though the ``correct''
answer is $\ell_c = 0$ (the subject's choice was predicted with certainty), the
floating-point arithmetic produces NaN, which then propagates through the
entire log-posterior sum.

\textbf{The fix.} Use $\epsilon = 10^{-7}$, which is above the float32 machine
epsilon. Then $1 - 10^{-7} \approx 0.9999999$ in float32, distinct from 1.0.
The clip now correctly prevents $\log(0)$, and $\log(10^{-7}) \approx -16.1$
contributes a finite (if extreme) value to the likelihood.

\textbf{Why the sampler didn't catch this.} The MCMC sampler evaluates the
log-posterior one block (one subject) at a time for Metropolis-within-Gibbs
proposals. A NaN log-posterior causes the MH step to reject the proposal
(NaN is not $>$ current), so the chain stays at the current state. The chain
never enters a NaN state. But the bridge sampling evaluator computes the
\textit{full joint} across all 111 subjects, and a single subject's NaN
poisons the sum. The sampler visits states where some subjects transiently
have extreme parameters that produce NaN in isolation, but those subjects
were updated at an earlier iteration and the chain has since moved on for
other subjects.

\subsection{When does the NaN occur?}

The NaN occurs when \textit{all} of the following hold simultaneously for
some subject $i$ in some chain state:

\begin{enumerate}
  \item The CRRA coefficient $r > 1$ (strong risk aversion), causing very flat
    normalized utility over the outcome range.
  \item The Prelec PWF parameters produce a decision weight $w \approx 1$
    for the preferred outcome.
  \item The noise parameter $\lambda$ is small ($\lambda < 0.01$), making
    $\Delta = (V_L - V_R) / (\lambda \cdot u_{\text{range}})$ very large.
  \item The observed choice matches the predicted preference, so $y_c = 1$
    and $P \approx 1$.
\end{enumerate}

These conditions occur for about 1--2 subjects per chain state in the
CRRA + Prelec model, affecting $\sim$48\% of saved chain states.


\section{The Proposal Distribution}

\subsection{Requirements}

The proposal $q(\bth)$ must satisfy:
\begin{enumerate}
  \item We can draw independent samples from $q$.
  \item We can evaluate $q(\bth)$ at any point (known normalizing constant).
  \item $q$ has sufficient overlap with the posterior: the ratio
    $\ptilde(\bth) / q(\bth)$ should not vary too wildly over the posterior
    support.
\end{enumerate}

\subsection{Multivariate normal}

The simplest choice: fit $q(\bth) = \mathcal{N}(\hat{\mu}, \hat{\Sigma})$
from the posterior sample mean and covariance. This satisfies requirements 1
and 2, but requirement 3 degrades rapidly with dimension. In $d = 452$
dimensions, even a well-fitted Gaussian generates samples that are far from
the posterior mode in at least some coordinates, producing extreme
$\ptilde / q$ ratios.

\subsection{Gaussian mixture}

A $K$-component Gaussian mixture with shared covariance:
\begin{equation}
  q(\bth) = \frac{1}{K} \sum_{k=1}^{K} \mathcal{N}(\bth \mid \mu_k, \Sigma)
\end{equation}
where $\mu_k$ are centers selected from the posterior samples (e.g., via
stratified sampling along the first principal component) and $\Sigma$ is the
full posterior covariance.

This improves coverage relative to a single Gaussian: proposal samples
cluster around actual posterior draws rather than all emanating from a single
center. Density evaluation is $O(K \cdot d^2)$ per point (one Mahalanobis
distance per component).

For the application in this project, $K = 10$ with the full $452 \times 452$
posterior covariance provides the best results.

\subsection{Kernel density estimate (KDE)}

Using $K = N_2$ (every posterior sample as a center) with a bandwidth
$h\Sigma$ is a Gaussian KDE. In principle this maximizes overlap, but in
$d = 452$ dimensions the Gaussian kernels are so concentrated that each
point only ``sees'' its nearest center. The effective density is approximately
constant across all points, providing no useful variation for the bridge
estimator.

\subsection{Factored proposals exploiting hierarchical structure}

The hierarchical model has conditional independence structure:
\begin{equation}
  p(\bth_1, \ldots, \bth_n, \bpsi \mid \by)
  = p(\bpsi \mid \by) \prod_{i=1}^{n} p(\bth_i \mid \bpsi, \by_i)
\end{equation}
A factored proposal $q(\bth, \bpsi) = q(\bpsi) \prod_i q(\bth_i \mid \bpsi)$
would only require fitting low-dimensional components (8-dimensional for
$\bpsi$, 3--4 dimensional for each $\bth_i$). However, this requires
modeling the dependence of $q(\bth_i \mid \bpsi)$ on $\bpsi$, which adds
implementation complexity.

This approach has not been tried in the current implementation but is the
most promising direction for improving the CRRA model's convergence.


\section{Convergence Diagnostics}

\subsection{What convergence means}

The iteration $\hat{r}^{(t+1)} = \ell_1^{(t)} / \ell_2^{(t)}$ should
converge to a fixed point. At each iteration, the change
$|\hat{r}^{(t+1)} - \hat{r}^{(t)}|$ should decrease geometrically toward
zero. A well-functioning estimator typically converges in 10--50 iterations.

\subsection{Monotone drift}

When the change decreases slowly or remains approximately constant, the
estimator is ``drifting.'' This occurs when the effective sample size is too
small relative to the dimension, typically because:
\begin{enumerate}
  \item A few samples with extreme $\ptilde/q$ ratios dominate one of the
    averages.
  \item The proposal has poor tail coverage, so posterior samples in the tails
    have very high $\ptilde/q$.
  \item Non-finite evaluations ($-\infty$ or NaN) create numerical artifacts
    in the log-sum-exp computation.
\end{enumerate}

\subsection{Initialization}

If $\hat{r}^{(0)}$ is far from the true $r$, the bridge function
degenerates. When $\hat{r}^{(0)} \ll r$, the term
$s_1 \ptilde$ dominates the denominator of $h^*$, and the estimator
reduces to the importance sampling estimator (which may have infinite
variance). When $\hat{r}^{(0)} \gg r$, the term $s_2 q r$ dominates,
and the estimator reduces to the harmonic mean estimator (which has
known instabilities).

A good initialization uses
$\log \hat{r}^{(0)} = \text{median}_{i}(\log \ptilde_i - \log q_i)$
over the posterior samples, which estimates $\log r$ without extreme values
dominating.


\section{Application: Hierarchical RDU Models}

\subsection{Model structure}

Each model has $n = 111$ subjects, each with $p = 3$ or $4$ structural
parameters, plus $H = 8$ hyperparameters. The full parameter vector has
dimension $d = n \cdot p + H = 341$ (Power) or $452$ (Prelec, GE).

The full unnormalized log-posterior is
\begin{equation}
  \log \ptilde = \underbrace{\log \pi(\bpsi)}_{\text{hyperprior}}
  + \sum_{i=1}^{111} \underbrace{\left[
    \log p(\bth_i \mid \bpsi) + \log |J_g(\bth_i)|
    + \log p(\by_i \mid \bth_i)
  \right]}_{\text{subject } i}
\end{equation}

Each subject contributes approximately $-50$ to the log-posterior (the data
likelihood for 80 binary choices). A ``bad'' evaluation from one subject
(e.g., $-500$ from NaN clamping, or $-\infty$ from overflow) dominates
the sum.

\subsection{Computational cost}

Evaluating $\log \ptilde$ at one point requires looping over 111 subjects,
each requiring utility computation, PWF evaluation, decision weights, RDU
values for 80 choices, and the Jacobian determinant. On GPU with JIT
compilation, this takes approximately 0.5--2ms per evaluation.

For bridge sampling with $N_1 = N_2 = 30{,}000$ samples, the total cost is
$60{,}000 \times 0.001 \approx 60$ seconds of GPU time, plus the overhead
of JIT compilation ($\sim$1--2 minutes) and density evaluation
($\sim$1 minute for the mixture proposal).

\subsection{Current results}

With the float32 clipping fix ($\epsilon = 10^{-7}$), a 10-component
Gaussian mixture proposal with full posterior covariance, and 2000 bridge
iterations:

\begin{center}
\begin{tabular}{lccc}
\toprule
Model & log ML & SE & Converged \\
\midrule
CRRA + Prelec & $\approx -5475$ & -- & No (drift $\sim$11/iter) \\
CRRA + GE & $\approx -5541$ & -- & Nearly (drift $\sim$0.0002/iter) \\
CRRA + Power & $-5614.22$ & 0.044 & Yes (iter 438) \\
\bottomrule
\end{tabular}
\end{center}

The ordering is CRRA $>$ GE $>$ Power (more complex models preferred),
with log Bayes factors of approximately $+77$ (CRRA vs GE) and $+73$
(GE vs Power). The data (8,880 binary choices) strongly support the
additional flexibility of the 2-parameter PWFs.

\subsection{Remaining issues}

The CRRA + Prelec model's bridge estimator has not converged despite
2000 iterations. The per-iteration drift of $\sim$11 log units suggests
poor effective overlap. The CRRA posterior spans a wider range of $r$ values
(including $r > 1$ where the likelihood has borderline numerical behavior),
and the 10-component mixture proposal does not adequately cover these
tails. Possible remedies:
\begin{enumerate}
  \item Increase the number of proposal samples ($N_1 \gg 30{,}000$).
  \item Use a factored proposal that exploits conditional independence.
  \item Switch to thermodynamic integration, which avoids the overlap problem.
\end{enumerate}


\section{Alternative Approaches}

\subsection{Thermodynamic integration (TI)}

TI estimates $\log p(\by \mid \cM)$ via
\begin{equation}
  \log p(\by \mid \cM) = \int_0^1 \E_\beta\!\left[\log p(\by \mid \bth)\right] d\beta
\end{equation}
where $\E_\beta$ is the expectation under the power posterior
$p_\beta(\bth) \propto p(\by \mid \bth)^\beta \pi(\bth)$.

TI requires running the sampler at $\sim$20--30 temperature levels
$\beta_k \in [0, 1]$ and recording the mean log-likelihood at each level.
The \texttt{bamcmc} sampler already supports parallel tempering, so the
infrastructure exists. TI avoids the proposal overlap problem entirely,
making it more robust in high dimensions, but requires substantially more
computation ($20\times$ the sampling cost).

\subsection{LOO-CV / WAIC}

Leave-one-out cross-validation via Pareto-smoothed importance sampling
\citep{Vehtari2017} requires only pointwise log-likelihood evaluations and
no proposal distribution. It estimates predictive performance rather than the
marginal likelihood, so it penalizes complexity differently from Bayes
factors. For the hierarchical models in this application, ``leave-one-out''
means leaving out one subject, and the LOO estimate requires the predictive
density $p(\by_i \mid \by_{-i})$ for each subject.

\subsection{Chib's method}

\citet{Chib1995} estimates the marginal likelihood from the posterior
ordinate at a single point $\bth^*$:
\begin{equation}
  p(\by \mid \cM) = \frac{p(\by \mid \bth^*) \pi(\bth^*)}{p(\bth^* \mid \by)}
\end{equation}
The numerator is easy to compute. The denominator $p(\bth^* \mid \by)$ is
estimated from the Gibbs full conditionals. For Metropolis-within-Gibbs,
this requires additional Rao-Blackwellization steps. The method is exact
for Gibbs samplers but introduces approximation error for MH-based samplers.


\bibliographystyle{apalike}
\begin{thebibliography}{99}

\bibitem[Chib, 1995]{Chib1995}
Chib, S. (1995).
\newblock Marginal likelihood from the {G}ibbs output.
\newblock {\em Journal of the American Statistical Association},
  90(432):1313--1321.

\bibitem[Gronau et~al., 2020]{Gronau2020}
Gronau, Q.~F., Sarafoglou, A., Matzke, D., Ly, A., Boehm, U.,
  Marsman, M., Leslie, D.~S., Forster, J.~J., Wagenmakers, E.-J., and
  Steingroever, H. (2020).
\newblock A tutorial on bridge sampling.
\newblock {\em Journal of Mathematical Psychology}, 94:102502.

\bibitem[Meng and Wong, 1996]{MengWong1996}
Meng, X.-L. and Wong, W.~H. (1996).
\newblock Simulating ratios of normalizing constants via a simple identity:
  A theoretical exploration.
\newblock {\em Statistica Sinica}, 6(4):831--860.

\bibitem[Vehtari et~al., 2017]{Vehtari2017}
Vehtari, A., Gelman, A., and Gabry, J. (2017).
\newblock Practical {B}ayesian model evaluation using leave-one-out
  cross-validation and {WAIC}.
\newblock {\em Statistics and Computing}, 27(5):1413--1432.

\end{thebibliography}

\end{document}
